\ea\label{ex:text15-1}
O paranda o peho taba vaa tea inu teara a ta Teapu.\\
\gll o	paranda	o	peho	taba	vaa	te=a	inu	te=ara	a	ta	teapu\\
     \textsc{art}3.\textsc{sg}	floor	\textsc{art}3.\textsc{sg}	\textsc{indef}3.\textsc{sg}	thing	\textsc{lnk}	\textsc{prep}=\textsc{art}2.\textsc{sg}	house	\textsc{prep}=1\textsc{pl}.\textsc{in}.\textsc{pron}	\textsc{art}2.\textsc{sg}	people	Teop\\
\glt `The floor is a thing in our house, the Teop people.'
\z

\ea\label{ex:text15-2}
Eove to kom, hio ae goroho raara.\\
\gll eove	to	kom	hio	ae	goroho	ra=ara\\
     \textsc{h}.3\textsc{sg}.\textsc{pron}	\textsc{rel}	step.on	sit	and1	sleep	1\textsc{pl}.\textsc{in}.\textsc{ipfv}=1\textsc{pl}.\textsc{in}.\textsc{pron}\\
\glt `It is where we step, sit, and sleep on.'
\z

\ea\label{ex:text15-3}
O paranda na vavaobete raara toon na naono vai to dao raara bono hogo.\\
\gll o	paranda	na	va-va-obete	ra=ara	toon	na	naono	vai	to	dao	ra=ara	bono	hogo\\
     \textsc{art}3.\textsc{sg}	floor	\textsc{tm}2	\textsc{mult}-\textsc{caus}-lie	1\textsc{pl}.\textsc{in}.\textsc{ipfv}=1\textsc{pl}.\textsc{in}.\textsc{pron}	top	3\textsc{sg}.\textsc{poss}=\textsc{art}3.\textsc{pl}	beam	\textsc{dem}6	\textsc{rel}	call	1\textsc{pl}.\textsc{in}.\textsc{ipfv}=1\textsc{pl}.\textsc{in}.\textsc{pron}	\textsc{obj}.\textsc{art}3.\textsc{sg}	floor.joist\\
\glt `We put the floor onto the beams that we call hogo (floor joists).'\footnote{lit. `The floor we put onto the beams ...'}
\z

\ea\label{ex:text15-4}
O paranda na pakupaku raara bono kobuu.\\
\gll o	paranda	na	paku{$\sim$}paku	ra=ara	bono	kobuu\\
     \textsc{art}3.\textsc{sg}	floor	\textsc{tm}2	\textsc{red}{$\sim$}make	1\textsc{pl}.\textsc{in}.\textsc{ipfv}=1\textsc{pl}.\textsc{in}.\textsc{pron}	\textsc{obj}.\textsc{art}3.\textsc{sg}	kobuu.palm\\
\glt `We make the floor from the \textit{kobuu} palm (a cyrtostachys species).'\footnote{lit. `The floor, we make from the \textit{kobuu} palm.'}
\z

\ea\label{ex:text15-5}
Tea meha maa vaan na pakupaku riori bono kumako.\\
\gll te=a	meha	maa	vaan	na	paku{$\sim$}paku	ri=ori	bono	kumako\\
     \textsc{prep}=\textsc{art}2.\textsc{sg}	other2.\textsc{sg}	\textsc{plm}	village	\textsc{tm}2	\textsc{red}{$\sim$}make	3\textsc{pl}.\textsc{ipfv}=3\textsc{pl}.\textsc{pron}	\textsc{obj}.\textsc{art}3.\textsc{sg}	kugumaru.palm\\
\glt `In other villages they make it from the kumako palm (areca novohibernica).'
\z

\ea\label{ex:text15-6}
Beara tau paku o paranda, eara toro ipinoo roho a bebeahu ae a tatabonata na inu bono kara urisi.\\
\gll be=ara	tau	paku	o	paranda	e-ara	toro	ipinoo	roho	a	bebeahu	ae	a	tatabonata	n=a	inu	bono	kara	urisi\\
     when=1\textsc{pl}.\textsc{in}.\textsc{pron}	\textsc{tm}5	make	\textsc{art}3.\textsc{sg}	floor	\textsc{h}-1\textsc{pl}.\textsc{in}.\textsc{pron}	\textsc{tm}7	measure	before	\textsc{art}2.\textsc{sg}	length	and1	\textsc{art}2.\textsc{sg}	width	3\textsc{sg}.\textsc{poss}=\textsc{art}3.\textsc{sg}	house	\textsc{obj}.\textsc{art}3.\textsc{sg}	string	wild.banana.sheath\\
\glt `When we are about to make the floor, we must measure the length and the width of the house with a tape made of the wild banana sheath.'
\z

\ea\label{ex:text15-7}
O kara urisi bona to kahi dee raara rasuu to kahi ipinoo raara o kobuu.\\
\gll o	kara	urisi	bona	to	kahi	dee	ra=ara	rasuu	to	kahi	ipinoo	ra=ara	o	kobuu\\
     \textsc{art}3.\textsc{sg}	string	wild.banana.sheath	\textsc{ana}.\textsc{sg}	\textsc{rel}	\textsc{tm}1	carry	1\textsc{pl}.\textsc{in}.\textsc{ipfv}=1\textsc{pl}.\textsc{in}.\textsc{pron}	jungle	\textsc{rel}	\textsc{tm}1	measure	1\textsc{pl}.\textsc{in}.\textsc{ipfv}=1\textsc{pl}.\textsc{in}.\textsc{pron}	\textsc{art}3.\textsc{sg}	kobuu.palm\\
\glt `This wild banana sheath tape we take to the jungle where we will measure the \textit{kobuu} palm.'
\z

\ea\label{ex:text15-8}
Beara suguna teo hum to tei nae o kobuu, eara paa sisira roho tea tabae o kobuu toro paku ara bono paranda toro vuaba vai o manoto.\\
\gll be=ara	suguna	te=o	hum	to	tei	nae	o	kobuu	e-ara	paa	siisira	roho	tea~tabae	o	kobuu	toro	paku	ara	bono	paranda	toro	vuaba	vai	o	manoto\\
     when=1\textsc{pl}.\textsc{in}.\textsc{pron}	arrive	\textsc{prep}=\textsc{art}3.\textsc{sg}	place	\textsc{rel}	be.there	3\textsc{sg}.\textsc{ipfv}	\textsc{art}3.\textsc{sg}	kobuu.palm	\textsc{h}-1\textsc{pl}.\textsc{in}.\textsc{pron}	\textsc{tm}3	examine	before	because	\textsc{art}3.\textsc{sg}	kobuu.palm	\textsc{tm}7	make	1\textsc{pl}.\textsc{in}.\textsc{pron}	\textsc{obj}.\textsc{art}3.\textsc{sg}	floor	\textsc{tm}7	one	\textsc{dem}6	\textsc{art}3.\textsc{sg}	mature\\
\glt `When we have arrived at the place where the \textit{kobuu} palm stands, we carefully examine (it), because the \textit{kobuu} palm from which we will make the floor must be a mature one.'
\z

\ea\label{ex:text15-9}
Beara viana vakavara o kobuu vai tau paku ara o paranda, eara repaa mosi gunaha e.\\
\gll be=ara	viana	va-kavara	o	kobuu	vai	tau	paku	ara	o	paranda	e-ara	re=paa	mosi	gunaha	e\\
     when=1\textsc{pl}.\textsc{in}.\textsc{pron}	choose	\textsc{advr}-complete	\textsc{art}3.\textsc{sg}	kobuu.palm	\textsc{dem}6	\textsc{tm}7	make	1\textsc{pl}.\textsc{in}.\textsc{pron}	\textsc{art}3.\textsc{sg}	floor	\textsc{h}-1\textsc{pl}.\textsc{in}.\textsc{pron}	\textsc{consec}=\textsc{tm}3	cut	go.down	3\textsc{sg}.\textsc{pron}\\
\glt `When we have chosen the \textit{kobuu} palm that we are about to make the floor from, we cut it down.'
\z

\ea\label{ex:text15-10}
O kobuu vai o kobusu saka mataa haana,\\
\gll o	kobuu	vai	o	kobusu	saka	mataa	haa=na\\
     \textsc{art}3.\textsc{sg}	kobuu.palm	\textsc{dem}6	\textsc{art}3.\textsc{sg}	moist	\textsc{neg}1	good	\textsc{neg}2=3\textsc{sg}.\textsc{ipfv}\\
\glt `The young and moist \textit{kobuu} palm is not good,'
\z

\ea\label{ex:text15-11}
eikou na antee vaha koana tea koava, evehee o vuaba vai o manoto vamataa na antee nana tea tei vananaoana.\\
\gll ei~kou	na	antee	vaha	koa=na	tea	koava	evehee	o	vuaba	vai	o	manoto	va-mataa	na	antee	nana	tea	tei	va-na{$\sim$}nao-ana\\
     because	\textsc{tm}2	can	again	just=3\textsc{sg}.\textsc{ipfv}	\textsc{compl}1	rot	but	\textsc{art}3.\textsc{sg}	one	\textsc{dem}6	\textsc{art}3.\textsc{sg}	mature	\textsc{advr}-good	\textsc{tm}2	can	3\textsc{sg}.\textsc{ipfv}	\textsc{compl}1	be.there	\textsc{advr}-\textsc{red}{$\sim$}go-\textsc{adjr}\\
\glt `because it can easily rot, but the fully mature one can last for a long time.'
\z

\ea\label{ex:text15-12}
Beo kobuu obete nana putaa, eara repaa ipinoo e,\\
\gll be=o	kobuu	obete	nana	putaa	e-ara	re=paa	ipinoo	e\\
     when=\textsc{art}3.\textsc{sg}	kobuu.palm	lie	3\textsc{sg}.\textsc{ipfv}	on.the.ground	\textsc{h}-1\textsc{pl}.\textsc{in}.\textsc{pron}	\textsc{consec}=\textsc{tm}3	measure	3\textsc{sg}.\textsc{pron}\\
\glt `When the \textit{kobuu} palm is lying on the ground, we measure it,'
\z

\ea\label{ex:text15-13}
eara repaa rom suku bata koara o ipinoo.\\
\gll e-ara	re=paa	rom	suku	bata	koa=ra	o	ipinoo\\
     \textsc{h}-1\textsc{pl}.\textsc{in}.\textsc{pron}	\textsc{consec}=\textsc{tm}3	cut.crosswise	follow	\textsc{cont}	just=1\textsc{pl}.\textsc{in}.\textsc{ipfv}	\textsc{art}3.\textsc{sg}	measure\\
\glt `and then we just cut the stem crosswise according to the measure.'
\z

\ea\label{ex:text15-14}
Beara romrom vakavara e, eara repaa paku pete a kiu vaa tea kopaa.\\
\gll be=ara	rom{$\sim$}rom	va-kavara	e	e-ara	re=paa	paku	pete	a	kiu	vaa	tea	kopaa\\
     when=1\textsc{pl}.\textsc{in}.\textsc{pron}	\textsc{red}{$\sim$}cut.crosswise	\textsc{advr}-complete	3\textsc{sg}.\textsc{pron}	\textsc{h}-1\textsc{pl}.\textsc{in}.\textsc{pron}	\textsc{consec}=\textsc{tm}3	do	subsequently	\textsc{art}2.\textsc{sg}	work	\textsc{lnk}	\textsc{compl}1	cut.lengthwise\\
\glt `When we have finished cutting it up crosswise, we subsequently do the work of cutting it lengthwise.'
\z

\ea\label{ex:text15-15}
Eara he toro kopaa e teo hum buaku.\\
\gll e-ara	he	toro	kopaa	e	te=o	hum	buaku\\
     \textsc{h}-1\textsc{pl}.\textsc{in}.\textsc{pron}	\textsc{advers}	\textsc{tm}7	cut.lengthwise	3\textsc{sg}.\textsc{pron}	\textsc{prep}=\textsc{art}3.\textsc{sg}	piece	two\\
\glt `But we must cut it lengthwise into two pieces.'
\z

\ea\label{ex:text15-16}
Tea kiu vaa tea kopaa, eara toro kiu nia buaku a toraaraa, ae amaa ta ginuva.\\
\gll te=a	kiu	vaa	tea	kopaa	e-ara	toro	kiu	ni=a	buaku	a	toraaraa	ae	a=maa	ta	ginuva\\
     \textsc{prep}=\textsc{art}2.\textsc{sg}	work	\textsc{lnk}	\textsc{compl}1	cut.lengthwise	\textsc{h}-1\textsc{pl}.\textsc{in}.\textsc{pron}	\textsc{tm}7	work	\textsc{appl}=\textsc{art}2.\textsc{sg}	two	\textsc{art}2.\textsc{sg}	axe	and1	\textsc{art}2.\textsc{sg}=\textsc{plm}	piece	wedge\\
\glt `For the work of the lengthwise cutting, we must work with two axes and some wedges.'
\z

\ea\label{ex:text15-17}
Tea kopaa, eara toro mosi momohu nia peha toraaraa,\\
\gll te=a	kopaa	e-ara	toro	mosi	momohu	ni=a	peha	toraaraa\\
     \textsc{prep}=\textsc{art}2.\textsc{sg}	cut.lengthwise	\textsc{h}-1\textsc{pl}.\textsc{in}.\textsc{pron}	\textsc{tm}7	cut	first	\textsc{appl}=\textsc{art}2.\textsc{sg}	one2.\textsc{sg}	axe\\
\glt `For the lengthwise cut, we must first make a cut with one axe'
\z

\ea\label{ex:text15-18}
ore topo vakisi nana, beara mosi nia vaabuaku.\\
\gll o=re	topo	vakisi	nana	be=ara	mosi	ni=a	vaa-buaku\\
     3.\textsc{pron}=\textsc{consec}	be.stuck	still	3\textsc{sg}.\textsc{ipfv}	when=1\textsc{pl}.\textsc{in}.\textsc{pron}	cut	\textsc{appl}=\textsc{art}2.\textsc{sg}	\textsc{ord}-two\\
\glt `so that it is still stuck (in the stem) when we make a cut with the second (axe).'
\z

\newpage
\ea\label{ex:text15-19}
A vaabuaku re topo vakisi nana, beara rahi pute a toraaraa momohu,\\
\gll a	vaa-buaku	re	topo	vakisi	nana	be=ara	rahi	pute	a	toraaraa	momohu\\
     \textsc{art}2.\textsc{sg}	\textsc{ord}-two	\textsc{consec}	be.stuck	still	3\textsc{sg}.\textsc{ipfv}	when=1\textsc{pl}.\textsc{in}.\textsc{pron}	pull	pull.out	\textsc{art}2.\textsc{sg}	axe	first\\
\glt `Then the second (axe) is still stuck (in the stem), when we pull out the first axe,'
\z

\ea\label{ex:text15-20}
beara tapa bata ni naora amaa ta ginuva teo hum to mosi vaakaa bata naona a bua toraaraa.\\
\gll be=ara	tapa	bata	ni	nao=ra	a=maa	ta	ginuva	te=o	hum	to	mosi	vaa-kaa	bata	nao=na	a	bua	toraaraa\\
     when=1\textsc{pl}.\textsc{in}.\textsc{pron}	drive	\textsc{cont}	\textsc{appl}	thither=1\textsc{pl}.\textsc{in}.\textsc{ipfv}	\textsc{art}2.\textsc{sg}=\textsc{plm}	piece	wedge	\textsc{prep}=\textsc{art}3.\textsc{sg}	place	\textsc{rel}	cut	\textsc{caus}-open	\textsc{cont}	thither=3\textsc{sg}.\textsc{ipfv}	\textsc{art}2.\textsc{sg}	two2.\textsc{sg}	axe\\
\glt `while we are driving the wedges into the stem where the two axes cut it lengthwise open.'
\z

\ea\label{ex:text15-21}
Eara pasi paku bata koa vo naora teebona, ore antee komana beo rom kobuu takopaa.\\
\gll e-ara	pasi	paku	bata	koa	vo	nao=ra	tee=bona	o=re	antee	komana	be=o	rom	kobuu	ta-kopaa\\
     \textsc{h}-1\textsc{pl}.\textsc{in}.\textsc{pron}	\textsc{tm}4	do	\textsc{cont}	just	like	thither=1\textsc{pl}.\textsc{in}.\textsc{ipfv}	\textsc{prep}=\textsc{ana}.\textsc{sg}	3.\textsc{pron}=\textsc{consec}	stop	altogether	when=\textsc{art}3.\textsc{sg}	stem	kobuu.palm	\textsc{acaus}-cut.lengthwise\\
\glt `We will do it like this, until the stem of the \textit{kobuu} palm is cut open lengthwise.'
\z

\ea\label{ex:text15-22}
Be kavara o hum kiu bona, eara repaa paku pete o hum vaa tea kopaa varutarutaa e.\\
\gll be	kavara	o	hum	kiu	bona	e-ara	re=paa	paku	pete	o	hum	vaa	tea	kopaa	va-ruta{$\sim$}rutaa	e\\
     when	finished	\textsc{art}3.\textsc{sg}	part	work	\textsc{ana}.\textsc{sg}	\textsc{h}-1\textsc{pl}.\textsc{in}.\textsc{pron}	\textsc{consec}=\textsc{tm}3	do	subsequently	\textsc{art}3.\textsc{sg}	part	\textsc{lnk}	\textsc{compl}1	cut.lengthwise	\textsc{advr}-\textsc{red}{$\sim$}small	3\textsc{sg}.\textsc{pron}\\
\glt `When this part of work is finished, we continue with the part of cutting it bit by bit lengthwise (into two halves).'
\z

\ea\label{ex:text15-23}
A kiu bona sa beera vaha haana,\\
\gll a	kiu	bona	sa	beera	vaha	haa=na\\
     \textsc{art}2.\textsc{sg}	work	\textsc{ana}.\textsc{sg}	\textsc{neg}1	big	again	\textsc{neg}2=3\textsc{sg}.\textsc{ipfv}\\
\glt `This work is not big again,'
\z

\ea\label{ex:text15-24}
na antee koana bona peha toraaraa.\\
\gll na	antee	koa=na	bona	peha	toraaraa\\
     \textsc{tm}2	need	only=3\textsc{sg}.\textsc{ipfv}	\textsc{art}2.\textsc{sg}	one2.\textsc{sg}	axe\\
\glt `(it) only needs one axe.'
\z

\ea\label{ex:text15-25}
Beara kopaa varutarutaa bata rae, eara mee toro ipinoo batara a tatabonata bona bebeahu na kobisi kuriraara.\\
\gll be=ara	kopaa	va-ruta{$\sim$}rutaa	bata	ra=e	e-ara	mee	toro	iponoo	bata=ra	a	tatabonata	bona	bebeahu	n=a	kubisi kuri-raara\\
     when=1\textsc{pl}.\textsc{in}.\textsc{pron}	cut.lengthwise	\textsc{advr}-\textsc{red}{$\sim$}small	\textsc{cont}	1\textsc{pl}.\textsc{in}.\textsc{ipfv}=3\textsc{sg}.\textsc{pron}	\textsc{h}-1\textsc{pl}.\textsc{in}.\textsc{pron}	also	\textsc{tm}7	measure	\textsc{cont}=1\textsc{pl}.\textsc{in}.\textsc{ipfv}	\textsc{art}2.\textsc{sg}	width	\textsc{obj}.\textsc{art}2.\textsc{sg}	length	3\textsc{sg}.\textsc{poss}=\textsc{art}2.\textsc{sg}	small.limb  hand-1\textsc{pl}.\textsc{in}.\textsc{poss}.\textsc{pron}\\
\glt `While we are cutting it open lengthwise bit by bit, we must also measure the width with the length of our finger.'
\z

\ea\label{ex:text15-26}
Be kavara a kiu bona, eara repaa tarataraha kahie bono ogana nae\\
\gll be	kavara	a	kiu	bona	e-ara	re=paa	tara{$\sim$}taraha	kahi=e	bono	ogana	nae\\
     when	finished	\textsc{art}2.\textsc{sg}	work	\textsc{ana}.\textsc{sg}	\textsc{h}-1\textsc{pl}.\textsc{in}.\textsc{pron}	\textsc{consec}=\textsc{tm}3	\textsc{red}{$\sim$}carve.off	from=3\textsc{sg}.\textsc{pron}	\textsc{obj}.\textsc{art}3.\textsc{sg}	soft.inside	3\textsc{sg}.\textsc{poss}.\textsc{pron}\\
\glt `When this work is finished, we carve its soft inside out'\footnote{lit. `we carve from it its inside'}
\z

\ea\label{ex:text15-27}
ae taraha taatagi koa amaa ta kobuu, ore paa kavara kou a kiu vaa teebona.\\
\gll ae	taraha	taatagi	koa	a=maa	ta	kobuu	o=re	paa	kavara	kou	a	kiu	vaa	tee-bona\\
     and1	carve.off	carefully	just	\textsc{art}2.\textsc{sg}=\textsc{plm}	piece	kobuu.palm	3.\textsc{pron}=\textsc{consec}	\textsc{tm}3	be.finished	\textsc{emph}2	\textsc{art}2.\textsc{sg}	work	\textsc{lnk}	\textsc{prep}-\textsc{ana}.\textsc{sg}\\
\glt `and just carefully carve off some pieces of kobuu palm, and then indeed the work on it is finished.'
\z

\ea\label{ex:text15-28}
A kiu vaa tea paku paranda a hihivaa,\\
\gll a	kiu	vaa	tea	paku	paranda	a	hihivaa\\
     \textsc{art}2.\textsc{sg}	work	\textsc{lnk}	\textsc{compl}1	make	floor	\textsc{art}2.\textsc{sg}	difficult\\
\glt `The work of making the floor is difficult,'
\z

\ea\label{ex:text15-29}
eikou o kobuu o taba kikis, ei mee vaaponis rakaha nana bona kiu bona.\\
\gll ei~kou	o	kobuu	o	taba	kikis	ei	mee	vaa-ponis	rakaha	nana	bona	kiu	bona\\
     because	\textsc{art}3.\textsc{sg}	kobuu.palm	\textsc{art}3.\textsc{sg}	thing	hard	\textsc{dem}2	also	\textsc{caus}-heavy	very	3\textsc{sg}.\textsc{ipfv}	\textsc{obj}.\textsc{art}2.\textsc{sg}	work	\textsc{ana}.\textsc{sg}\\
\glt `because the kobuu palm is a hard thing, and this also makes this work very heavy.'
\z

\ea\label{ex:text15-30}
Ei, eara deedee voosu vai, are paa no vaobete e tea inu, ore paa tavus paranda.\\
\gll ei	e-ara	dee{$\sim$}dee	voosu	vai	a=re	paa	no	va-obete	e	te=a	inu	o=re	paa	tavus	paranda\\
     \textsc{dem}2	\textsc{h}-1\textsc{pl}.\textsc{in}.\textsc{pron}	\textsc{red}{$\sim$}carry	come.home	now	1\textsc{pl}.\textsc{in}.\textsc{pron}=\textsc{consec}	\textsc{tm}3	go	\textsc{caus}-lie	3\textsc{sg}.\textsc{pron}	\textsc{prep}=\textsc{art}2.\textsc{sg}	house	3.\textsc{pron}=\textsc{consec}	\textsc{tm}3	become	floor\\
\glt `Now, we carry it home, then we make it lie in the house, so that it becomes a floor.'
\z
